\section{Zero Knowledge TIPA Scheme}
\label{sec:zkgipa}
We derive from the GIPA (General Inner Pairing Argument) \cite{Bunz21}
and design the scheme for a lookup argument for group elements.
The original GIPA scheme is not zero knowledge. This section provides
an zero knowledge extension of it, inspired by \cite[Appendix E]{Bunz21}.
The basic idea is to pad the commitment with random nonces and
cancel them in pairing equations.

The GIPA scheme relies on an adaption of
AFGHO16 commitment \cite{Abe16} to commit to a vector of group elements.
Given a bilinear pairing group $(\F,\G_1, \G_2, \G_T, e)$, and an integer
$n$. The prover key of AFGHO commitment is a sequence
$\left(\alpha^{i}]_1, [\beta^{i}]_2\right)_{i=0}^{2n-2}$ for trapdoor
$\alpha$ and $\beta$. Then given a vector of group elements
$\myvec{u} \in \G_1^n$, its AFGHO commitment
$\gcm_1(\myvec{u}) = \prod_{i=1}^n e(\myvec{u}_i, [\beta^{2(i-1)}]_2)$.
Similaly for $\myvec{v} \in \G_2^n$, 
$\gcm_2(\myvec{v}) = \prod_{i=1}^n e([\alpha^{2(i-1)}]_1, \myvec{v}_i)$.
Notice that the above commitments are not hiding, but
it can be fixed by adding random nonce.
As AFGHO commitment generates elements in $\G_T$, we sometimes write
it as $\GCommit_{\myvec{v}}$ using \ttt{mathscr} font, to distinguish
from commitments in $\G_1$ or $\G_2$ (e.g., $\Commit_{\myvec{u}}$ the
Pedersen vector commitment to vector $\myvec{u}$).
Another annotation difference is that $\G_T$ is written as a multiplicative
group but $\G_1$ is written as additive group. For instance,
$e([a+b]_1, [c]_2) = e([a]_1, [c]_2) e([b]_1 [c]_2)$
and $e(a[v]_1, [1]_2) = e([v]_1, [1]_2)^a$.

We now formally define the GIPA scheme below and we'll call them
as blackbox functions in our own protocol. It is proved in \cite{Bunz21}
that the GIPA scheme is knowledge sound (i.e., the prover can prove
its knowledge of the vector of source vector of group and field elements).


\begin{definition} \cite{Bunz21}
    Let $B=(\F,\G_1,\G_2,\G_T,e)$ be a bilinear group for a given security
parameter $\lambda$. The generalized inner pairing product argument (GIPA)
provides the following algorithms: 
\begin{enumerate}
\item
$(\pk_G, \vk_G) = \Setup(\powparam,B,n)$, where: \newline
$\pk_G = \left(\left([\alpha^{2(i-1)}]_1, 
[\beta^{2(i-1)}]_2\right)_{i=1}^{2n}, [\mu]_1, [\mu]_2, [\nu]_1,
\right)$ for some trapdoor $\alpha$, $\beta$, 
$\mu$ and $\nu$.
Let $\vk_G$ $=$ $([\alpha]_2,$ $[\beta]_1$, $[\mu]_2, [\nu]_1)$.
\item 
$\h_1 \leftarrow \gcm_1(\myvec{a}, \pk_G)$: \newline
It generates the commitment
for a group element vector $\myvec{a} \in \G_1^{n}$.
\footnote{Note that it is required that $|\myvec{a}|=n+1$, and
this is verified by the well-formed check step in GIPA.}
Similarly are $\gcm_2$ and $\gcm_{\F}$ defined for input vectors
of $\G_2^n$ and $\F^n$.
When the context is clear we omit the commitment key
$\pk_G$ in notations, e.g., $\gcm_1(\myvec{a})$.

\item $(Z, \prf) \leftarrow \prove(\myvec{a}, \myvec{b}, r, \pk_G)$: \newline
It generates the inner product argument $Z$ and the proof $\prf$.
GIPA has three instantiations: $\tipp$,
$\mippu$, and $\mippk$: 
    $\tipp$: $\myvec{a} \in \G_1^{n}$, $\myvec{b} \in \G_2^{n}$,
        $\Commit_{\myvec{a}} = \gcm_1(\myvec{a})$,
        $\Commit_{\myvec{b}} = \gcm_2(\myvec{b})$,
        $Z = \sum_{i=1}^{n}
            \left(r^{2(i-1)} (e(\myvec{a}_i,\myvec{b}_i))\right)$.
    $\mippu$: $\myvec{a} \in \G_1^{n}$, $\myvec{b} \in \F^{n}$,
        $\Commit_{\myvec{a}} = \gcm_1(\myvec{a})$,
        $\Commit_{\myvec{b}} = \gcm_{\F}(\myvec{b})$,
        $Z = \sum_{i=1}^{n} \left(r^{2(i-1)} \myvec{b}_i \myvec{a}_i \right)$.
        $\mippk$ differs from $\mippu$ in that
        $\Commit_{\myvec{b}} = \myvec{b}$.

\item $1/0 \leftarrow \verify(\Commit_{\myvec{a}}, \Commit_{\myvec{b}},$ $r, Z, \prf, \vk_G)$: verifies the proof.

\end{enumerate}
\end{definition}

We now construct a zero knowledge version of the GIPA scheme,
and focus on TIPP only.

\begin{definition}\label{zk_tipp}
 Given $\myvec{a} \in \G_1^n$ and $\myvec{b} \in \G_2^n$, the
 $\zktipp$ argument provides the following operation:

\begin{enumerate}
\item
 $(\GCommit_\myvec{a}, \GCommit_\myvec{b}, \Z, \Commit_{r_z}, \prf)$
   $\leftarrow$ $\zktippprove(\myvec{a}, \myvec{b}, r):$ \newline
   It samples two random nonces $r_1, r_2 \in \F$, and use them to
   provide hiding AFGHO commitments of the two vectors,
   and a proof for $\Z = \prod_{i=1}^n e(r^{2(i-1)}\myvec{a}_i, \myvec{b}_i)
   ~ e([r_z]_1, [1]_2)$, which is $r_z$ away from the real
   inner pairing product and $r_z$ hides in Pedersen commitment
   $\Commit_{r_z} = [r_z]_1 + r_z'[\nu]_2$.  Since $r_z$ is random,
   the scheme does not leak information of $\myvec{a}$ and $\myvec{b}$
   and their inner pairing product.

\item
 $1/0 \leftarrow \zktippverify(\GCommit_{\myvec{a}}, \GCommit_{\myvec{b}}, r, \Z, \Commit_{r_z}, \prf)$: \newline
 Given two hiding AFGHO commitments, verify their well-formedness (i.e.,
 commitment to $n$ elements), and the validity of the proof.
\end{enumerate}
\end{definition}




We now present the details of the $\zktipp$ protocol in
Figure\ \ref{fig_zktipp}.

\def\hp{\hspace*{0.1in}}
\begin{small}
\begin{algorithm}[bt]
\SetKwData{Left}{left}\SetKwData{This}{this}\SetKwData{Up}{up}
\SetKwFunction{Union}{Union}\SetKwFunction{FindCompress}{FindCompress}
$ (\GCommit_\myvec{a}, \GCommit_\myvec{b}, \Z, \Commit_{r_z}, \prf) 
	\leftarrow \newline
	\zktippprove(\myvec{a} \in \G_1^n, r_a \in \F, \myvec{b} \in \G_2^n, 
		r_b \in \F, r \in \F, z \in \F, r_z \in \F, \pk, \vk)$
\newline
(1) Parse $\pk = \left(\left([\alpha^i]_1, [\beta^i]_2\right)_{i=0}^{2n-2},
	[\mu]_1, [\mu]_2, [\nu]_1\right)$. 
Sample $r_4, r_5, r_6, r_7, r_8$ from $\F$, 
and sample $\myvec{r_a}, \myvec{r_b}$ from $\F^n$. Then compute:
\[
\begin{array} {l}
\GCommit_\myvec{a} \leftarrow \gcm_1(\myvec{a})~ e([[r_a]_1, [\mu]_2) \\
\GCommit_\myvec{b} \leftarrow \gcm_2(\myvec{b})~ e([r_b]_1, [\mu]_2) \\
\Z \leftarrow \prod_{i=1}^n e(r^{2(i-1)} \myvec{a}_i, \myvec{b}_i)
~ e(z [1]_1, [1]_2) \\
\Commit_{z} \leftarrow [z]_1 + r_z[\nu]_1 \\
\GCommit_\myvec{r_a} \leftarrow \gcm_1(\myvec{r_a})~ e([r_4]_1, [\mu]_2) \\
\GCommit_\myvec{r_b} \leftarrow \gcm_2(\myvec{r_b})~ e([r_5]_1, [\mu]_2)\\
\newline
\Commit_\myvec{r_a\times b} \leftarrow 
\sum_{i=1}^n (r^{2(i-1)} [\myvec{r_a}_i) \myvec{b}_i]_1 + 
  r_6[\mu]_1 \\
\Commit_\myvec{r_b\times a} \leftarrow 
\sum_{i=1}^n (r^{2(i-1)} \myvec{r_b}_i) \myvec{a}_i + 
  r_7 [\mu]_1 \\
\Commit_\myvec{r_b\times r_a} \leftarrow 
\sum_{i=1}^n r^{2(i-1)}[\myvec{r_a}_i \myvec{r_b}_i]_1
 + r_8 [\mu]_1
\end{array}
\]
 (2) Fiat-Shamir: $c_1, c_2 \leftarrow$
 	$\hash\big(\vk,\GCommit_\myvec{a}, \GCommit_\myvec{b}, r, \Z,$
		$\Commit_{r_z}, $
		$\GCommit_\myvec{r_a}, \GCommit_\myvec{r_b},$
		$\Commit_\myvec{r_a\times b},$
		$\Commit_\myvec{r_b\times a},$
		$\Commit_\myvec{r_a\times r_b}\big)$. \newline
 (3) Compute: \newline
	\[
		\begin{array}{l}
 \myvec{A} \leftarrow c_1 \myvec{a} + [\myvec{r_a}]_1  \\
 \myvec{B} \leftarrow c_2 \myvec{b} + [\myvec{r_b}]_2  \\
 \GCommit_{\myvec{A}} \leftarrow \gcm_1(\myvec{A})  \\
 \GCommit_{\myvec{B}} \leftarrow \gcm_2(\myvec{B})  \\
 \Z_2 \leftarrow \prod_{i=1}^n e(r^{2(i-1)} [\myvec{A}_i]_1,
		[\myvec{B}_i]_2)	  \\
 \prf_t \leftarrow \tipp.\prove(\myvec{A}, \myvec{B}, r, \Z_2) \\
 	   \rho_1 \leftarrow -(c_1 r_a + r_4)  \\
	   \rho_2 = -(c_2 r_b + r_5) \\
 	\B \leftarrow 
 	e(-r_z [c_1 c_2 \nu]_1 + r_6 [c_2 \mu]_1
 	+ r_7 [c_1 \mu]_1]  + r_8 [\mu]_1, [1]_2)  \\
	\prfdlog_{\B\_r_z} \leftarrow \prfdlog(\B,\Commit_{r_z}) 
	  \big\{(r_z, r_z', r_k): \\
	  \hspace*{0.1in}\B =  
	  	 e([-c_1c_2\nu]_1, [1]_2)^{r_z'} ~e([\mu]_1,  [1]_2) ^{r_k} \\
	  \hspace*{0.1in}\Commit_{r_z} = r_z[1]_1 + r_z'[\nu]_2 \\
	  \big\}
   \end{array} 
	\]
 (4) Let $\prf = (\GCommit_\myvec{r_a}, $
 			$\GCommit_\myvec{r_b}, $
 			$\Commit_\myvec{r_a \times b}, $
 			$\Commit_\myvec{r_b \times a}, $
 			$\Commit_\myvec{r_a \times r_b}, $
 			$\Commit_\myvec{r_a \times r_b}, $
 			$\rho_1, \rho_2, $
			$\GCommit_\myvec{A}$, $\GCommit_\myvec{B}$,
			$\Z_2$,
			$\prf_t$, $\B$,
			$\prfdlog_{\B\_r_z}$.
		Return $(\GCommit_{\myvec{a}}, \GCommit_{\myvec{b}}, \Z, \Commit_{r_z},
				\prf)$

$1/0 \leftarrow \zktippverify(\vk, 
	\Commit_{\myvec{a}}, 
	\Commit_{\myvec{b}}, r, \Z, \Commit_{r_z}, \prf):$  \newline
	Parse $\prf$ as shown in $\prove$ and returns $1$
	if and only if the following pass (given $c_1, c_2$ computed
	using $\prf$ as shown in $\zktippprove$).  \newline
	(1) $\verify(\prfdlog_{\B\_{r_z}})$.  \newline
	(2) $\GCommit_\myvec{A} =  \GCommit_\myvec{a}^{c_1}  ~\GCommit_\myvec{r_a}
		 e([\rho_1], [\mu]_2)$. \newline
	(3) $\GCommit_\myvec{B} =  \GCommit_\myvec{b}^{c_2}  ~\GCommit_\myvec{r_b}
		 e([\rho_2], [\mu]_2)$. \newline
	(4) $\Z_2\B = 
			(\Z / e(\Commit_{r_z}, [1]_2))^{c_1 c_2} ~
			e(, \Commit_{r_b \times a}, [1]_2)^{c_1} \newline
			\hspace*{0.1in} e([1]_2, \Commit_{r_a \times b})^{c_2}~
			e(\Commit_{r_a \times r_b}, [1]_2)$.\newline
	(5) $\tipp.\verify(\GCommit_\myvec{A}, \GCommit_\myvec{B}, r, \Z_2)$.
\caption{Zero Knowledge TIPP Protocol}
\label{fig_zktipp}
\end{algorithm}
\end{small}

\begin{lemma}\label{lemma_zktipp} Under the $q$-ASDBP and $q$-SDH assumptions
\cite{Bunz21} in AGM and ROM, the construction provided in 
Figure\ \ref{fig_zktipp} is perfectly complete, computational knowledge
sound and perfectly zero knowledge.
\end{lemma}
\begin{proof} The $q$-ASDBP and $q$-SDH assumptions are needed by the
GIPA scheme \cite{Bunz21}. The completeness is apparent. We first
prove the knowledge soundness.  In the following we assume
all notations of $\prf$ in the $\zktippverify()$ operation.

Compared with the standard TIPP scheme \cite{Bunz21}, 
our prover key has extra keys: $([\nu]_1, [\mu]_2, [\mu]_1)$. 
Thus, a group element provided by the prover
in $\G_1$ (and other groups similarly)
can be written as a polynomial $[f(\alpha, \beta, \mu, \nu)]_1$
in the AGM model. We call $f$ the {\em representing polynomial}
(over variables $\alpha$, $\beta$, $\mu$, $\nu$)
of the group element. When a polynomial is defined over
$\alpha$ and $\beta$ only, we call it {\em standard} (i.e.,
corresponding to those in the standard TIPP scheme, because
$\mu$ and $\nu$ are the extra variables we have introduced).
We also call the group element {\em standard} if its reprenting
polynomial is standard.

From Check (5), based on the knowledge soundness of GIPA scheme,
the prover has the knowledge of $\myvec{A} \in \G_1^n$ 
and $\myvec{B} \in \G_2^n$ s.t.
$\GCommit_a = \gcm_1(\myvec{A})$.
$\GCommit_b = \gcm_2(\myvec{B})$.
Note that $\myvec{A}$ and $\myvec{B}$ are
standard (i.e., its representing polynomial
is defined over variables $\alpha$ and $\beta$ only).
Then from Check (2) and (3) we can tell that
there exist 
$\myvec{a}, \myvec{r_a} \in \G_1^n$ and
$\myvec{b}, \myvec{r_b} \in \G_2^n$
s.t. $\myvec{A} = c_1 \myvec{a}  + \myvec{r_a}$
and $\myvec{B} = c_2 \myvec{b} + \myvec{r_b}$,
and all of them are standard, because they are all committed
before verifier challenge $c_1$ and $c_2$ are raised and the
random challege force them to be standard because $\myvec{A}$
and $\myvec{B}$ are standard.
It then follows from (5) that:
\begin{small}
\begin{equation} \label{eq_z2}
  \Z_2 = \prod_{i=1}^n \left( 
  	e(\myvec{a}_i, \myvec{b}_i)^{c_1c_2}
  	e(\myvec{a}_i, [\myvec{r_b}_i]_2)^{c_1}
  	e(\myvec{[r_a}_i]_1, \myvec{b}_i)^{c_2}
  	e(\myvec{[r_a}_i]_1, [\myvec{r_b}_i]_2)
		\right)^{r^{2(i-1)}}
\end{equation}
\end{small}

From Check (1), the DLOG proof provides the knowledge
of two commitments, so that we know that there exist
$r_z, r_z', r_{u}, r_{v}, r_{w}$
s.t. the following is true. Note that 
$r_z'$ appears in both $\B$ and $\Commit_{r_z}$,
which is used later in Equation\ \ref{eq_c1c2}
to cancel terms. $r_{u}$, $r_v$ and $r_w$, with
the information of $c_1$ and $c_2$ can be extracted from
the value of $r_2$. As $\mu$ and $\nu$ are generated randomly (independent
of each other), $\B$ can be regarded as a Pedersen commitment
which has perfect hiding and computational binding (i.e.,
it is hard for an adversary to fake value of $r_z'$).
\begin{small}
\begin{eqnarray}
 \B &=& 
 	e([c_1c_2\nu]_1, [1]_2)^{r_z'} ~
 	e([c_1\mu]_1, [1]_2)^{r_u} \nonumber \\
	& &
 	e([c_2\mu]_1, [1]_2)^{r_v}~
 	e([\mu]_1, [1]_2)^{r_w} \label{eq_B} \\
 \Commit_{r_z} &=& [r_z]_1 + r_z'[\mu] \label{eq_rz}
\end{eqnarray}
\end{small}

Consider $\Commit_{r_a \times b}$, it is not guaranteed
to be standard. Write it as
$\Commit_{r_a \times b} = \Commit_{r_a \times b, 1}
 + \Commit_{r_a \times b, \mu}$, where
 $\Commit_{r_a \times b, 1}$ corresponds to part
 defined over $\alpha$ and $\beta$ (including $[1]_1$
 and $[1]_2$), and
 $\Commit_{r_a \times b, \mu}$ for the rest.
Apply similar notation to
$\Commit_{r_b \times a}$ and $\Commit_{r_a \times r_a}$.
Integrating Check (4) with
Equations \ref{eq_z2}, \ref{eq_B}, \ref{eq_rz}, we have:

\begin{small}
\begin{eqnarray} 
 e([0]_1,[0]_1]) &=& \left(\frac{\prod_{i=1}^n r^{2(i-1)}
 		e(\myvec{a}_i, \myvec{b}_i)}
			{Z - e([1]_1, [1]_1)^{r_z} }\right)^{c_1 c_2} 
				\label{eq_c1c2} \\
	& & e\left(\Commit_{\myvec{r_b} \times \myvec{a},1}, [1]_2\right)^{c1} \label{eq_c1} \\
	& & e\left([1]_1, \Commit_{\myvec{r_a} \times b,1}\right)^{c2} \label{eq_c2} \\
	& & e\left([1]_1, \Commit_{\myvec{r_a} \times \myvec{r_b},1}\right) \label{eq_c3} \\
	& &
		\left(e(\Commit_{\myvec{r_a} \times \myvec{b},\mu}, [1]_2) / e([\mu]_1, [1])^{r_u}\right)^{c1} \label{eq_c1mu} \\
	& &
		\left(e([1]_1, \Commit_{\myvec{r_b} \times \myvec{a},\mu}) / e([\mu]_1, [1])^{r_v}\right)^{c2} \label{eq_c2mu} \\
	& &
		\left(e(\Commit_{\myvec{r_a} \times \myvec{r_b},\mu}, [1]_2) / e([\mu]_1, [1])^{r_w}\right) \label{eq_c3mu}
\end{eqnarray}
\end{small}

% Consider the sum of Lines \ref{eq_c1c2} to \ref{eq_c3}, 
% its AGM representation can be regarded as one multivariate
% polynomial over $\alpha, \beta, c_1, c_2$ (not including $\mu$),
% where the co-efficients are already fixed (i.e., $r$, $\myvec{a}_i$,
% $\myvec{b}_i$, $r_z$ and the AGM representation of $\Z$).
% 
% For Lines \ref{eq_c1mu} to \ref{eq_c3mu}, its representative
% polynomial (let it be $P(\alpha, \beta, \mu, c_1, c_2)$
% are further parameterized by $r_u$, $r_v$, and $r_w$
% controlled by a potentially malicious prover. We first argue that
% $P(.)$ must be a zero polynomial, otherwise,
% based on how $\Commit_{\myvec{r_a} \times \myvec{b}}$ and etc.
% are split, once can infer that each non-zero monomial
% of $P(.)$ must contains $\mu$ (which is unknown to the prover)
% it yields a non-zero value with high probability
Consider the AGM representation polynomial for the
RHS of Lines \ref{eq_c1c2} to \ref{eq_c3mu}.
It must be a zero polynomial.  Given each line represents
a sub-polynomial that are grouped by distinct set of
variables that it contain (e.g., $c_1 c_2$ for 
Line \ref{eq_c1c2}), each sub-polynomial should be
a zero polynomial as well. This leads to
the main result:
$\Z = \sum_{i=1}^n e(\myvec{a}_i, \myvec{b_i}) + e([r_z]_1, [1])_2)$
where $r_z$ hides in $\Commit_{r_z}$.

{\bf Zero Knowledge:} the proof is straightforward. The simulator
samples $\myvec{A}$, $\myvec{B}$, which uniquely determines
$\GCommit_\myvec{A}$, $\GCommit_\myvec{B}$, and $\Z_2$, and
$\prf_t$.
Then the simulator samples $\rho_1$, $\rho_2$, 
$c_1$, $c_2$, $\GCommit_\myvec{a}$,
$\GCommit_\myvec{b}$, 
$\Commit_{\myvec{r_a} \times \myvec{b}}$,
$\Commit_{\myvec{r_b} \times \myvec{a}}$,
$\Commit_{\myvec{r_a} \times \myvec{r_b}}$. Note that they
are independent of each other (following the same distribution of
a real conversation).
Then the above uniquely determines: 

\begin{footnotesize}
\begin{eqnarray*}
& & \Commit_{\myvec{r_a}} \leftarrow 
\frac{\GCommit_{\myvec{A}} }
{ \GCommit_{\myvec{a}}^{c_1} e([\rho_1]_1, [\mu]_2) } \\
& & \Commit_{\myvec{r_b}} \leftarrow 
\frac{\GCommit_{\myvec{B}} }
{ \GCommit_{\myvec{B}}^{c_1} e([\rho_2]_1, [\mu]_2) } \\
& &\B \leftarrow  \\
& &\frac{
			(\Z / e(\Commit_{r_z}, [1]_2))^{c_1 c_2}
			e(\Commit_{r_b \times a}, [1]_2)^{c_1}~
			e([1]_2, \Commit_{r_a \times b})^{c_2}~
			e(\Commit_{r_a \times r_b}, [1]_2)}
			{ \Z_2 }
\end{eqnarray*}
\end{footnotesize}

Finally, run the simulator for DLOG of
$\Commit_{r_z}$ and $\Commit_{\B}$ for 
$\prfdlog_{\B\_r_z}$.
\end{proof}

We have a similar zero knowledge construction for 
MIPP, as shown in Figure\ \ref{fig_zkmipp}.
In this case $\myvec{a} \in \G_1^n$
and $\myvec{b} \in \F^n$.
In this case $Z = \sum_{i=1}^n r^{2(i-1)} \myvec{a_i}^{b_i} + r_z [1]_1$,
and $\Commit_{r_z} = r_z[1]_1 + r_3[\mu]_1$, and both of them
are in $\G_1$. The details are presented in Figure\ \ref{fig_zkmipp}.


\begin{lemma}\label{lemma_zkmipp} Under the $q$-ASDBP and $q$-SDH assumptions
\cite{Bunz21} in AGM and ROM, the construction provided in 
Figure\ \ref{fig_zkmipp} is perfectly complete, computational knowledge
sound and perfectly zero knowledge.
\end{lemma}

\def\hp{\hspace*{0.1in}}
\begin{small}
\begin{algorithm}[bt]
\SetKwData{Left}{left}\SetKwData{This}{this}\SetKwData{Up}{up}
\SetKwFunction{Union}{Union}\SetKwFunction{FindCompress}{FindCompress}
$ (\GCommit_\myvec{a}, \Commit_\myvec{b}, Z, \Commit_{r_z}, \prf) 
	\leftarrow \newline
	\zkmippprove(\myvec{a} \in \G_1^n, \myvec{b} \in \F^n, r \in \F, \pk, \vk)$
\newline
(1) Parse $\pk = \left(\left([\alpha^i]_1, [\beta^i]_2\right)_{i=0}^{2n-2},
	[\mu]_1, [\mu]_2, [\nu]_1\right)$. 
Sample $r_1, r_2, r_z, r_3, r_4, r_5, r_6, r_7, r_8$ from $\F$, 
and sample $\myvec{r_a}, \myvec{r_b}$ from $\F^n$. Then compute:
\[
\begin{array} {l}
\GCommit_\myvec{a} \leftarrow \gcm_1(\myvec{a})~ e([r_1]_1, [\mu]_2) \\
\Commit_\myvec{b} \leftarrow \gcm_\F(\myvec{b})~ r_2 [\mu]_2 \\
\Z \leftarrow \prod_{i=1}^n r^{2(i-1)} \myvec{b}_i \myvec{a}_i
+ r_z [1]_1 \\
\Commit_{r_z} \leftarrow [r_z]_1 + r_3[\nu]_1 \\
\GCommit_\myvec{r_a} \leftarrow \gcm_1(\myvec{r_a})~ e([r_4]_1, [\mu]_2) \\
\Commit_\myvec{r_b} \leftarrow \gcm_\F(\myvec{r_b}) + r_5 [\mu]_1\\
\newline
\Commit_\myvec{r_a\times b} \leftarrow 
\sum_{i=1}^n (r^{2(i-1)} \myvec{r_a}_i) \myvec{b}_i + 
  r_6[\mu]_2 \\
\Commit_\myvec{r_b\times a} \leftarrow 
\sum_{i=1}^n (r^{2(i-1)} \myvec{r_b}_i) \myvec{a}_i + 
  r_7 [\mu]_1 \\
\Commit_\myvec{r_b\times r_a} \leftarrow 
\sum_{i=1}^n r^{2(i-1)}[\myvec{r_a}_i \myvec{r_b}_i]_1
 + r_8 [\mu]_1
\end{array}
\]
 (2) Fiat-Shamir: $c_1, c_2 \leftarrow$
 	$\hash\big(\vk,\GCommit_\myvec{a}, \Commit_\myvec{b}, r, \Z,$
		$\Commit_{r_z}, $
		$\GCommit_\myvec{r_a}, \Commit_\myvec{r_b},$
		$\Commit_\myvec{r_a\times b},$
		$\Commit_\myvec{r_b\times a},$
		$\Commit_\myvec{r_a\times r_b}\big)$. \newline
 (3) Compute: \newline
	\[
		\begin{array}{l}
 \myvec{A} \leftarrow c_1 \myvec{a} + [\myvec{r_a}]_1  \\
 \myvec{B} \leftarrow c_2 \myvec{b} + [\myvec{r_b}]_2  \\
 \GCommit_{\myvec{A}} \leftarrow \gcm_1(\myvec{A})  \\
 \Commit_{\myvec{B}} \leftarrow \gcm_\F(\myvec{B})  \\
 Z_2 \leftarrow \sum_{i=1}^n r^{2(i-1)} \myvec{B}_i [\myvec{A}_i]_1, \\
 \prf_t \leftarrow \mipp.\prove(\myvec{A}, \myvec{B}, r, Z_2) \\
 	   \rho_1 \leftarrow -(c_1 r_1 + r_4)  \\
	   \rho_2 \leftarrow -(c_2 r_2 + r_5) \\
 	\B \leftarrow 
 	e(-r_3 [c_1 c_2 \nu]_1 + r_6 [c_2 \mu]_1
 	+ r_7 [c_1 \mu]_1]  + r_8 [\mu]_1, [1]_2)  \\
	\prfdlog_{\B\_r_z} \leftarrow \prfdlog(\B,\Commit_{r_z}) 
	  \big\{(r_z, r_z', r_k): \\
	  \hspace*{0.1in}\B =  
	  	 e([-c_1c_2\nu]_1, [1]_2)^{r_z'} ~e([\mu]_1,  [1]_2) ^{r_k} \\
	  \hspace*{0.1in}\Commit_{r_z} = r_z[1]_1 + r_z'[\nu]_2 \\
	  \big\}
   \end{array} 
	\]
 (4) Let $\prf = (\GCommit_\myvec{r_a}, $
 			$\Commit_\myvec{r_b}, $
 			$\Commit_\myvec{r_a \times b}, $
 			$\Commit_\myvec{r_b \times a}, $
 			$\Commit_\myvec{r_a \times r_b}, $
 			$\Commit_\myvec{r_a \times r_b}, $
 			$\rho_1, \rho_2, $
			$\GCommit_\myvec{A}$, $\Commit_\myvec{B}$,
			$Z_2$,
			$\prf_t$, $\B$,
			$\prfdlog_{\B\_r_z}$.
		Return $(\GCommit_{\myvec{a}}, \Commit_{\myvec{b}}, \Z, \Commit_{r_z},
				\prf)$

$1/0 \leftarrow \zkmippverify(\vk, 
	\GCommit_{\myvec{a}}, 
	\Commit_{\myvec{b}}, r, Z, \Commit_{r_z}, \prf):$  \newline
	Parse $\prf$ as shown in $\prove$ and returns $1$
	if and only if the following pass (given $c_1, c_2$ computed
	using $\prf$ as shown in $\zktippprove$).  \newline
	(1) $\verify(\prfdlog_{\B\_{r_z}})$.  \newline
	(2) $\GCommit_\myvec{A} =  \GCommit_\myvec{a}^{c_1}  ~\GCommit_\myvec{r_a}
		 e([\rho_1], [\mu]_2)$. \newline
	(3) $\Commit_\myvec{B} =  {c_2} \Commit_\myvec{b} + \Commit_\myvec{r_b} +
		 \rho_2 [\mu]_1$. \newline
	(4) $e(Z_2, [1]_2) \B =$ \newline 
			\hspace*{0.1in} $e(c_1c_2(Z - \Commit_{r_z}) +
			c_1\Commit_{r_b \times a} +
			c_2\Commit_{r_a \times b} + 
			\Commit_{r_a \times r_b}, 
			[1]_2)$ \newline
			\newline
	(5) $\mipp.\verify(\GCommit_\myvec{A}, \Commit_\myvec{B}, r, Z_2)$.
\caption{Zero Knowledge MIPP Protocol}
\label{fig_zkmipp}
\end{algorithm}
\end{small}

\section{ZK-VPD Scheme}
\label{sec:zkvpd}
Our scheme depends on the univariate zk-VPD scheme
\cite{Zkreg}, which is an instantiation of zero knowledge 
polynomial delegation scheme for univariate polynomials \cite{zkvsql}.
It provides a zero knowledge enhancement
of the KZG polynomial commitment scheme. For completeness of the paper,
we present a slighly modified scheme from what is presented in
\cite{Zkreg,zkvsql}, by allowing its set-up to take
randoms sampled by the caller procedure. The details
of the protocol 
are presented in Fig\ \ref{alg:kzg}. We refer readers to
\cite{Zkreg,zkvsql} for the details of the proof of the soundness and
zero knowledge properties.

Intuitively, as shown in Figure\ \ref{alg:kzg}, the $\CommitPoly()$
function generates a hiding KZG commitment to a polynomial, and
its knowledge proof (which simulates the KZG knowledge proof
that leverages the knowledge of exponent assumption). 
The $\CommitValue()$ function generates a proof that
a hiding polynomial evaluates to a hidden value $y$
at a public point $t$, but $y$ hides in a Pedersen
commitment $\Commit_y$. Thus, it does not no leak information
of the committed polynomial.

\def\hp{\hspace*{0.1in}}
\begin{scriptsize}
\begin{algorithm}[bt]
\SetKwData{Left}{left}\SetKwData{This}{this}\SetKwData{Up}{up}
\SetKwFunction{Union}{Union}\SetKwFunction{FindCompress}{FindCompress}
$(\pk, \vk) \leftarrow \vpdsetup(\alpha, \nu, n)$
\newline
Sample $\beta$ from $\F$.
Compute $\gamma \leftarrow \left(
\Bigl((\alpha)^i\Bigr)^{n}_{i=0},
\Bigl(\beta(\alpha)^i\Bigr)^{n}_{i=0},
\nu,
\beta \nu,
%\Bigl((s_2)^i\Bigr)^{q}_{i=0},
%\Bigl(\alpha(s_2)^i\Bigr)^{q}_{i=0},
\alpha \nu
%, \alpha s_1s_2
\right)$.
\newline
Return $(\pk = [\gamma]_1, \vk=([\alpha]_2, [\beta]_2, [\alpha \beta]_2, [\alpha \nu]_2))$.
\newline
\\
($\Commit_{p},\prf_p) \leftarrow
	 \CommitPoly(p(X), r_p, \pk)$:  \hp \newline
Return $\left([p(\alpha) + r_p \nu]_1, [\beta(p(\alpha)  + r_p \nu)]_1\right)$.
\newline
\\
$1/0 \leftarrow \CheckCommit(\Commit_p,\prf_{p},\vk)$:
 \hp Return 1 if and only if $e(\Commit_p, [\beta]_2) = e(\prf_{p},
	[1]_2)$.
\newline
\\
$(\Commit_y,\prf) \leftarrow
	\CommitValue(p(X), r_p, r_y, t, \pk)$:  \newline
Compute: $y \leftarrow p(t)$.
$q_1(X) \leftarrow (p(X) - y)/(X-t)$.  \newline
Sample $r_q$ from $\F$.
Compute: \newline
$\Commit_y \leftarrow [y]_1 + r_y[\nu]_1$. \newline
$\Commit_1 \leftarrow [q_1(\alpha)]_1 + r_q [\nu]_1$. \newline 
$\Commit_2 \leftarrow r_q[\nu(\alpha-t)]_1 + (r_y-r_p)[\nu]_1$. \newline
Compute Schnorr proofs: \newline
$\prfdlp_{y} \leftarrow \prfdlp(\Commit_y, ([1]_1, [\nu]_1))
		\{(y,r_y): \Commit_y = y[1]_1 + r_y [\nu]_1\}$, and
$\prfdlp_{2} \leftarrow \prfdlp(\Commit_2, ([\nu(\alpha-t)]_1, [\nu]_1))\{(r_1, r_2): \Commit_2 = r_1[\nu(\alpha-t)]_1 + r_2[\nu]_1\}$.
\newline
Let $\prf = (\Commit_1, \Commit_2, \prfdlp_{y}, \prfdlp_{2})$.
Return $(\Commit_y, \prf)$.
\newline

$1/0 \leftarrow \CheckValue(\Commit_p,\Commit_y,t,\prf,\vk)$:  \newline
{\bf Assumption:}  $\exists \prf_p$ for $\Commit_p$
s.t. $\CheckCommit(\Commit_p, \prf_p, \vk)=1$.
\newline
Parse $\prf$ as $(\Commit_1, \Commit_2, \prfdlp_{y}, \prfdlp_{2})$.
Return 1 if and only if: \newline
$
\CheckDLP(\Commit_y,\prfdlp_{y}, ([1]_1, [\nu]_1))=1 \And \CheckDLP(\Commit_2,\prfdlp_{2}, ([\nu(\alpha-t)]_1, [\nu]_1))=1
\And e\left(\Commit_p - \Commit_y + \Commit_2\right), [1]_2)
= e(\Commit_1, [\alpha- t]_2).
$


\caption{Univariate zk-VPD Scheme}
\label{alg:kzg}
\end{algorithm}
\end{scriptsize}

Porting the proof in \cite{Zkreg} to AGM, we have the following.

\begin{lemma} Under the QDLOG assuption in AGM, the scheme presented
in Figure\ \ref{alg:kzg} is perfectly complete, computational
knowledge sound and binding, and perfectly zero knowledge.
\end{lemma}

\section{Argument for Equivalence Between AFGHO commitment and Pedersen
Vector Commitment}
\label{sec:eqcom}
This section provides a construction that asserts  
the equivalence of various forms of commitment
to a vector $\myvec{x} \in \F^n$:
\begin{enumerate}
	\item AFGHO16 commitment: $\GCommit_{\myvec{x}} = \gcm_1([\myvec{x}]_1) + e(r_1[\mu]_1, [1]_2)$. When the context is clear (i.e., it is known that
	$[\mu]_1$ is used as the extra key for blinding factor),
	we write it as $\GCommit_{\myvec{x}} = \gcm_1([\myvec{x}]_1; r_1)$.
	\item Pedersen vector commitment over Lagrange basis:
	$\Commit_{\myvec{x}} = \sum_{i=1}^n \myvec{x}_i 
		[L_{n,i}(s)]_1 + r_2[\mu]_1$. 
	\item KZG commitment to $\myvec{x}$ as co-efficients:
		$\KCommit_{\myvec{x}} = \sum_{i=1}^n \myvec{x}_i [\alpha^{i-1}]_1 + r_3[\mu]_1$.
\end{enumerate}

The construction is formally defined below:

\begin{enumerate}
	\item $(\pk,\vk) \leftarrow \comeqsetup(n, s):$ 
	Generate the prover and verifker keys for vector of size $n$.
	$s$ is the trapdoor information given by the caller (as a part
	of a larger setup process).

	\item $(\GCommit_{\myvec{x}}, \Commit_{\myvec{x}}, \KCommit_{\myvec{x}}, \prf) \leftarrow \comeqprove(\myvec{x}, r_1, r_2, r_3, \pk)$: 
	Generate commitments to $\myvec{x}$ in different forms using
	the randoms $r_1$ to $r_3$. Provide a proof which asserts
	their equivalence.

	\item $1/0 \leftarrow \comeqverify(\GCommit_{\myvec{x}}, 
		\Commit_{\myvec{x}}, \KCommit_{\myvec{x}}, \prf, \vk)$.
\end{enumerate}
\def\hp{\hspace*{0.1in}}

Although it might be possible to split the construction into
several pair-wise equivalence proof, we pack them into one
for the convenience of the presentation. We given an intuitive
explaination of the idea below.
We first leverage the QA-NIZK (used in 
LegoSnark \cite{LegoSnarkFull}) to prove the equivalence
of a Pedersen Vector Commitment (over Lagrange basis) and
a KZG commitment. Then we use the zk-MIPP argument to prove
the equivalence of the AFGHO and the KZG, which leverages
the zk-VPD scheme as shown in Section\ \ref{sec:zkvpd}. 
The construction is shown in Figure\ \ref{fig_eqv}.


\begin{footnotesize}
\begin{algorithm}[bt]
\SetKwData{Left}{left}\SetKwData{This}{this}\SetKwData{Up}{up}
\SetKwFunction{Union}{Union}\SetKwFunction{FindCompress}{FindCompress}
$ (\pk,\vk) \leftarrow \comeqsetup(n, s)$
\newline
(1) Compute $(\pk_{\ttt{G}}, \vk_{\ttt{G}}) 
	\leftarrow {\texttt{ZkGIPA.\texttt setup}}(n)$,
	Parse and retrive $([\alpha]_1)_{i=1}^n$ and $[\mu]_1$
	from $\pk_{\ttt{G}}$ and $\vk_{\ttt{G}}$ of zk-GIPA.
	Define $\Commit_{\ttt{1}} = \gcm_\F\left(\left([1]_1)_{i=1}^n\right)\right)$.\newline
	Compute QA-NIZK keys for KZG and Pedersen vector commitments:
	\newline
    $\myvec{M} \leftarrow \left(\begin{array}{l l}
		{[}L_{n,1}(s)]_1, ..., [L_{n,n}(s)]_1, & [\mu]_1, [0]_1 \\
		{[}\alpha^0]_1, ..., [\alpha^{n-1}]_1, & [0]_1, [\mu]_1 
		\end{array}
	\right)$. \newline
	Sample $\myvec{k} \in \F^2$, $a\in \F$. Compute
	$\myvec{P} = \myvec{M}^\intercal \myvec{k} \in \G_1^{n+1}$,
	$\myvec{C} = a \cdot [\myvec{k}]_2 \in \G_2^2$. \newline
(2) $(\pk_{\ttt{V}}, \vk_{\ttt{V}}) \leftarrow
		\vpdsetup(\alpha, \mu, n)$. \newline
(3) Return $(\pk=(\pk_{\ttt{G}}, \pk_{\ttt{V}}, \myvec{M}, \myvec{P}), 
	\vk=(\vk_{\ttt{G}}, \vk_{\ttt{V}}, \myvec{C}, [a]_2, \Commit_{\ttt{1}})$. \newline

$(\GCommit_{\myvec{x}}, \Commit_{\myvec{x}}, \KCommit_{\myvec{x}}, \prf)
	\leftarrow \comeqprove(\pk, \myvec{x}, r_1, r_2, r_3)$:  \newline
(1) Compute 
	$\GCommit_{\myvec{x}} \leftarrow \gcm_1(\myvec{x}; r_1)$.  ~~
	$\myvec{w} \leftarrow \myvec{x} \concat r_2 \concat r_3$.  ~~
	$(\Commit_{\myvec{x}}, \KCommit_{\myvec{x}}) 
		\leftarrow \myvec{M} \myvec{w}$. \newline
(2) Fiat-Shamir: $r \leftarrow \hash(\vk, \GCommit_{\myvec{x}}, \Commit_{\myvec{x}}, \KCommit_{\myvec{x}})$. \newline
(3) Define $p(X) = \sum_{i=1}^n \myvec{x}_i X^{i-1}$.
	Define $y = \sum_{i=1}^{n} r^{2(i-1)} \myvec{x}_i$.
	Sample $r_y$ from $\F$ and compute: 
	$(\Commit_y, \prf_y) = \CommitValue(p(X), r_3, r_y, r^2, \pk_{\ttt{V}})$.
	\newline
(4) Compute using $r_1$ for $\GCommit_{x}$ in:
	$(\GCommit_{x}, \Commit_{\ttt{one}}, Z, \Commit_{r_z}, \prf_{\ttt{m}})
	\leftarrow \zkmippprove(\myvec{x}, ([1]_1)_{i=1}^n, r, \pk_{\ttt{G}},
	\vk_{\ttt{G}})$. \newline
	Let $Z = [y+z]_1$ and $\Commit_{r_z} = [z]_1 + r_z[\mu]_1$ in the
	above. \newline
(5) Let $D = \Commit_{r_z} + \Commit_{y} - Z$. 
	Compute: \newline $\prf_{\ttt{D1}} 
		\leftarrow \prfdlog(D,\Commit_{\ttt{one}};$ 
	$[\mu]_1)$$\{(u,v,w):$ $D = u[\mu]_1 + v[\nu]_1,$ $\Commit_{\ttt{one}}-\Commit_{\ttt{1}}=w[\mu]_1\}$.
	Note that for honest prover: $u = r_y$ and $v=r_z$. \newline
(6)  $\prf_{\ttt{qa}} \leftarrow \myvec{w}^\intercal \myvec{P}$. \newline
(7) Return $(\GCommit_{\myvec{x}}, \Commit_{\myvec{x}}, \KCommit_{\myvec{x}},
	\prf=(\prf_{\ttt{D1}}, \Commit_y, \prf_y, \Commit_{\ttt{one}}, 
		Z, \Commit_{r_z},
		\prf_{\ttt{qa}},\prf_{\ttt{m}}))$. \newline

$1/0 \leftarrow \comeqverify(\vk, 
	\GCommit_{\myvec{x}}, \Commit_{\myvec{x}}, \KCommit_{\myvec{x}},
	\prf):$  \newline
Parse $\prf$ as in $\comeqprove$, compute $r$ using Fiat-Shamir,
and return $1$ if and only if all of the following checks pass. \newline
(1) $e(\prf_{\ttt{qa}}, [a]_2) = e(\Commit_{\myvec{x}}, \myvec{C}_1) 
	~e(\KCommit_{\myvec{x}}, \myvec{C}_2)$. \newline
(2) $\zkmippverify(\GCommit_{x}, r, Z, \Commit_{r_z}, \prf_{\ttt{m}})$. \newline
(3) $\CheckValue(\KCommit_{\myvec{x}}, \Commit_y, r^2, \prf_{y}, \vk_{\ttt{G}})$. \newline
(4) Let $D = \Commit_{r_z} + \Commit_y - Z$.
   $\verify(\prf_{\ttt{D1}})$.

\caption{Equivalence Between Hiding AFGHO Commitment and Pedersen Vector Commitment}
\label{fig_eqv}
\end{algorithm}
\end{footnotesize}

\begin{lemma}\label{lemma_comeq} Under the $q$-ASDBP and 
$q$-SDH assumptions \cite{Bunz21} in AGM and ROM, the construction provided in 
Figure\ \ref{fig_gl} is perfectly complete, computational knowledge
sound and perfectly zero knowledge.
\end{lemma}
\begin{proof} We focus on knowledge soundness and zero knowledge.
The prover's knowledge of $[\myvec{x}]_1$, and $\myvec{x}$ behind 
$\GCommit_\myvec{x}$, $\Commit_{\myvec{x}}$, and $\KCommit_{\myvec{x}}$
are established using the zk-MIPP proof, and the QA-NIZK proof.
Check (1) in $\comeqverify()$
establishes the equivalence between $\Commit_{\myvec{x}}$
and $\KCommit_{\myvec{x}}$. Since Check (4) estalibhses that
$\Commit_{\ttt{one}}$ is a hiding commitment to $([1]_1)_{i=1}^n$,
from Check (2), one can infer that the prover knows
$z, r_5, r_6$ s.t.
\begin{eqnarray}
 & & Z = [z]_1 + [r_5] \label{eqZ} \\
 & & \Commit_{r_z} = [r_5]_1 + r_6 [\mu]_1 \label{eq_rz}
\end{eqnarray}

Then from Check (3) we have: the prover has knowledge of
$y$ and $r_7$ s.t.
\begin{equation} \label{eq_y}
\Commit_y = [y]_1 + r_y[\mu]_1
\end{equation}

Combining Equations (\ref{eqZ},\ref{eq_rz},\ref{eq_y}), $D$ can be written as
\begin{equation} \label{eq_d1}
D = (y-z)[1]_1 + (r_y + r_6) [\mu]_1
\end{equation}

From Check (4): the prover knows $r_8$ s.t.
\begin{equation} \label{eq_d2}
D = r_8 [\mu]_1
\end{equation}

Combining Equations \ref{eq_d1} and \ref{eq_d2}, we have:
\[
	(y-z)[1]_1 = (r_8-r_y-r_6) [\mu]_1
\]
If $y\neq z$ and $r_8\neq r_y + r_6$, the prover can compute
the dislogarithm and retrieve trapdoor $\mu$. Then we have
$y=z$, combining with check (2) and (3) again, and based 
on that $r$ is the verifier challenge after $\KCommit_{\myvec{x}}$
and $\GCommit_{\myvec{x}}$ are sent, we reach
the equivalene between
$\GCommit_{\myvec{x}}$ and $\KCommit_{\myvec{x}}$.

{\bf Zero Knowledge:} $\Commit_y$,
$\Commit_{\ttt{one}}$, $Z$, $\Commit_{r_z}$ in the proof are
randomly sampled by the simulator, thus 
independent of each other (and the claim), and the rest:
$\prf_{\ttt{D1}}$, $\prf_{\ttt{qa}}$, $\prf_{\ttt{m}}$
follow the same distribution (as provided) by the corresponding
simulator of the component proof.


\end{proof}

\section{Position Hiding Lookup Arguments for Group Elements}
Baloo \cite{Baloo} provides a position hiding look-up argument
for vector of field elements. We consider its variation for 
group elements: Given a public look up table
$\myvec{T} \in \G_1^N$, and a hiding AFGHO commitment 
$\gcm_1(\myvec{t})$ to 
a query table $\myvec{t} \in \G_1^n$, and a Pedersen vector commitment
$\Commit_{\myvec{x}}$ to a vector $\myvec{x} \in \F^n$, show that
$\myvec{t} \lookup \myvec{T}$, in addition:
for each $i\in [n]: \myvec{t}_{i} = \myvec{T}_{\myvec{x}_i}$.
The prover cost is $O(n)$ pairings and verifier cost is
$O(\log(n))$ $\G_T$ operations. Our solution is 
{\em trapdoor} in that the trusted set-up is ``per circuit" and it knows
the discrete logarithm of each $\myvec{T}_i$ (i.e.,
it knows the $u \in \F$ s.t. $\myvec{T}_i = [u]_1$, and the
prover does not have to know it). 

The construction is suitable for performing the ``selection"
of verifier keys of a collection of Groth'16 zkSNARK circuits,
thus the trusted setup of Groth'16 scheme can be used to
set up the presented lookup argument for the public lookup table
of group elements .  We name the scheme $\zkglook$. 
The design is based on zero knowledge GIPA scheme introduced
in Section \ref{sec:zkgipa}, which extends the GIPA scheme
in \cite{Bunz21}. The formal definition of the construction
is provided below:

\begin{enumerate}
	\item $([\myvec{T}]_1,\pk,\vk) \leftarrow \glsetup(\myvec{T}\in \F^N):$ 
	Assume the security parameter and bilinear groups are already
	given.  A trusted set-up, given a vector of field elements 
	$\myvec{T} \in \F^N$, generates the lookup table $[\myvec{T}]_1
	\in \G_1^N$, and the corresponding prover and verifier keys.


	\item $(\GCommit_{\myvec{t}}, \Commit_{\myvec{x}}, \prf)
	\leftarrow \glprove(\myvec{t} \in \G_1^n, r_t \in \F, \myvec{x}\in \F^n, 
		r_x \in \F, \pk)$: 
	It produces a zero knowledge proof $\prf$ for
	$\myvec{t} \lookup [\myvec{T}]_1$, and for
	each $i\in [n]$:
	$\myvec{t}_i = \myvec{T}_{\myvec{x}_i}$.
	Hiding Perdesen vector commitment
	to $\myvec{t}$ and $\myvec{x}$ are generated with the proof.

	\item $1/0 \leftarrow \glverify(\GCommit_{\myvec{t}}, 
		\Commit_{\myvec{x}}, \prf, \vk):$ 
	It checks that $\prf$ is valid.
\end{enumerate}
\def\hp{\hspace*{0.1in}}


\begin{footnotesize}
\begin{algorithm}[bt]
\SetKwData{Left}{left}\SetKwData{This}{this}\SetKwData{Up}{up}
\SetKwFunction{Union}{Union}\SetKwFunction{FindCompress}{FindCompress}
$ ([\myvec{T}]_1, \pk,\vk) \leftarrow \glsetup(\myvec{T} \in \F^N, n)$
\newline
(1) 
Define ${T}(X) = \sum_{i=1}^N \myvec{T}_i \cdot L_{N,i}(X)$,
Sample $s\in \F$. \newline
Compute: $[T(s)]_1$ and 
 	for $i\in [N]:$ $\prf_i \leftarrow
	\left[\frac{T(s) - \myvec{T}_i}{s-\omega_{N}^i}\right]_2$. \newline
(2) Let $\myvec{1}_n = (1)_{i=1}^n$.
Compute:
 \begin{eqnarray*}
  & & (\pk_{\ttt{eq}}, \vk_{\ttt{eq}}) \leftarrow
	\comeqsetup(n). \\
  & &
	(\pk_{\ttt{cq}}, \vk_{\ttt{cq}}) \leftarrow
	\cqsetup(s, N, n). \\
  & &
	\GCommit_{\ttt{1}} \leftarrow \gcm_2\left([\myvec{1_n}]_2 \right)
 \end{eqnarray*}
(3) Return 
$\big(\pk = \left(\pk_{\ttt{eq}}, \pk_{\ttt{cq}}, (\prf_i)_{i=1}^N,  \right)$, 
$\vk = \left(\vk_{\ttt{eq}}, \vk_{\ttt{cq}}, \GCommit_{T}, \GCommit_{s}, \GCommit_{1} \right) \big)$. \newline


$(\GCommit_{\myvec{t}}, \Commit_{\myvec{x}}, \prf)
	\leftarrow \glprove(\myvec{t}, r_t, \myvec{x}, r_x, \pk, \vk)$: \newline
(1) Parse $\pk$ and $\vk$ as shown in $\texttt{setup}$.  \newline
(2) Sample $r_{x_1}, r_{x_3}, z_u, r_{z_u}, z_v, r_{z_v}, r_u, r_v, r_1, 
r_{\pi}$ 
from $\F$ and compute: \newline
\begin{eqnarray*}
& &
	(\GCommit_{\myvec{x}}, \Commit_{\myvec{x}}, \KCommit_{\myvec{x}},
	\prf_{\ttt{eq}})
	\leftarrow \comeqprove(\myvec{x}, r_{x_1}, r_{x}, r_{x_2}, r_{x_3}, 
		\pk_{\ttt{eq}}) \\
& & \GCommit_{\myvec{t}} = \gcm_1(\myvec{t}) ~e(r_t[\mu]_1, [1]_2) \\
& & (\Commit_{\myvec{x}'}, \prf_{\ttt{eq}}) \leftarrow \cqprove(\myvec{x}, r_{x_2},\pk_{\ttt{cq}}) 
\end{eqnarray*}
(3) Fiat-Shamir: $r \leftarrow \hash(\vk, \GCommit_{\myvec{t}}, \Commit_{\myvec{x}})$. \newline
(4) Define $\myvec{u} \in \G_1^n$ s.t.
	$\myvec{u} = \left([T(s)]_1\right)_{i=1}^n - \myvec{t}$.
 	$\myvec{v} = ([s]_1)_{i=1}^n - [\myvec{x}]_1$.  
	Retrieve zk-GIPA keys $\pk_{\ttt{G}}$, 
	$\vk_{\ttt{G}}$ from $\pk_{\ttt{eq}}$.
Compute the following using $\pk_{\ttt{G}}$, $\vk_{\ttt{G}}$:
\begin{eqnarray*}
& &
(\GCommit_{\myvec{u}}, \GCommit_{\ttt{one}}, \Z_u, \Commit_{\Z_u},
	\prf_u) \leftarrow
	\zktippprove(\myvec{u}, r_u, [\myvec{1}_n]_2, r_1, r, z_u, r_{z_u})\\
& &
( \GCommit_{\myvec{v}}, \GCommit_{\myvec{\prf}}, \Z_v, \Commit_{\Z_v}, \prf_\myvec{v}) \leftarrow
	\zktippprove(\myvec{v}, r_v, (\prf_i)_{i=1}^n, r_{\pi}, r, z_v, r_{z_v})
	\\
& &
\prf_{\ttt{DL}} = \prfdlog(\Commit_{\Z_u}, \Commit_{\Z_v}, 
	\GCommit_{\tt{one}}, \GCommit_{u}, \GCommit_{T}, \GCommit{t},
	\GCommit{v}, \GCommit{s}, \GCommit{x}, \Commit_{\myvec{x}}, \Commit_{\myvec{x}'}) \\
& &\{(h,j,u,v,w,x,y): \\ 
& & \hspace*{0.1in}
	\Commit_{{\Z_u}} - \Commit_{{\Z_v}} =
		x[1]_1 + y[\nu]_1 \\
& & \hspace*{0.1in}
	\Commit_{\myvec{x}'} - \Commit_{\myvec{x}} =  h [z_\V(s)]_1 + j [\mu]_1 \\
& & \hspace*{0.1in}
	\Z_u/\Z_v = e([1], [1])^x \And  \\
& & \hspace*{0.1in}
	\GCommit_{\myvec{u}} \GCommit_{\myvec{t}} / \GCommit_{T}  
		=  e([\mu]_1, [1]_2)^u \And  \\
& & \hspace*{0.1in}
	\GCommit_{\myvec{v}} \GCommit_{\myvec{x}} / \GCommit_{s}  
		=  e([\mu]_1, [1]_2)^v \And  \\
& & \hspace*{0.1in}
	\GCommit_{\tt{one}} / \GCommit_{1} =  e([\mu]_1, [1]_2)^w \And  \\
& &\}
\end{eqnarray*}
where $h=r_{x_2}$, $j=-r_x$, $x = z_u - z_v$, $y = r_{z_u} - r_{z_v}$, $u = r_u + r_t$,
$v = r_v + r_{x_1}$, $w = r_1$. \newline
(4) Let $\prf = ($
	$\GCommit_{\myvec{x}}, \KCommit_{\myvec{x}},$
	$\prf_{\ttt{eq}}, $
	$\prf_{\ttt{cq}}, $
	$\GCommit_{\myvec{u}}, $
	$\GCommit_{\ttt{one}},$ 
	$\Z_u,$ $\Commit_{\Z_u},$ $\prf_u,$
	$\GCommit_{\myvec{v}}, $
	$\GCommit_{\myvec{\prf}},$
	$\Z_v, $
	$\Commit_{\Z_v},$ 
	$\prf_v,$ 
	$\prf_{\ttt{DL}})$.
	Return $(\GCommit_{\myvec{t}}, \Commit_{\myvec{x}}, \prf)$. \newline

$1/0 \leftarrow \glverify(\vk, 
	\GCommit_{\myvec{t}}, \Commit_{\myvec{x}}, \prf):$  \newline
Parse $\prf$ as shown in $\prove$ and returns $1$
	and compute $r$ using Fiat-Shamir. \newline
(1) $\zktippverify(\GCommit_{\myvec{u}}, \GCommit_{\ttt{one}}, r, \Z_u, 
	\Commit_{\Z_u}, \prf_u)$. \newline
(2) $\zktippverify(\GCommit_{\myvec{v}}, \GCommit_{\myvec{\prf}}, 
	r, \Z_v, \Commit_{\Z_v}, \prf_v)$. \newline
(3)  $\verify(\prf_{\ttt{DL}}; \Commit_{\Z_u}, \Commit_{\Z_v}, 
	\GCommit_{\tt{one}}, \GCommit_{u}, \GCommit_{T}, \GCommit{t},
	\GCommit{v}, \GCommit{s}, \GCommit{x}, \Commit_{\myvec{x}}, \Commit_{\myvec{x}'})$  \newline
(4) $\cqverify(\Commit_{\myvec{x}}, \vk_{\ttt{cq}})$ and
$\comeqverify(\GCommit_{\myvec{x}}, \Commit_{\myvec{x}}, \KCommit_{\myvec{x}},
\vk_{\ttt{eq}})$.

\caption{Position Hiding Group Elements Lookup Protocol} 
\label{fig_gl}
\end{algorithm}
\end{footnotesize}

We provide the intuition of the algorithm in Figure\ \ref{fig_gl} below.
Essentially, we are proving the following equation for
all $i\in [n]:$
\[
	e([T(s)]_1 - \myvec{t}_i) = e([s-\myvec{x}_i], \prf_i)
\]
This can be essentially achieved using the zero knowledge GIPA scheme,
following the stragegy of aggregating Groth16 proofs introduced in
\cite{Bunz21}. We then need to justify that each $\myvec{x}_i$ is a
valid index in range $[0,n]$. This is accomplished using
zk-$\cq$. To connect the $\GCommit_{\myvec{x}}$ used in zk-GIPA scheme
and the $\Commit_{\myvec{x}}$ used in zk-$\cq$, the
equivalence of commitments proofs introduced in \ref{sec:eqcom}
is used.

\begin{lemma}\label{lemma_gl} Under the $q$-ASDBP and 
$q$-SDH assumptions \cite{Bunz21} in AGM and ROM, the construction provided in 
Figure\ \ref{fig_gl} is perfectly complete, computational knowledge
sound and perfectly zero knowledge.
\end{lemma}

\begin{proof} (Sketch) The proof techniques are similar to previous
constructions in the appendix and we provide a sketch here. 

For knowledge soundness: Check (4) establishes that all elements in
$\myvec{x}$ are valid indices. $\prf_{\ttt{DL}}$ in Check (3) shows 
$\GCommit_{\ttt{one}}$ is a hiding commitment to $([1]_1)_{i=1}^n$. 
Then, the zk-GIPA verify functions
in (1) and (2) estalishes the knowledge of 
$\GCommit_{\myvec{u}} = \gcm_1(\myvec{u}) + r_u e([1]_1, [1]_2)$,
and similarly $\GCommit_{\myvec{v}}$. Then, in Check (3),
the DLOG proof on 
$\GCommit_{\myvec{u}}\GCommit_{\myvec{t}}/\GCommit_{\myvec{T}}$
and
$\GCommit_{\myvec{v}}\GCommit_{\myvec{x}}/\GCommit_{x}$
show the prover's knowledge of $\myvec{t}$ and $\myvec{x}$
behind $\GCommit_{\myvec{t}}$ and $\GCommit_{\myvec{x}}$. 
They also prove the following:
\begin{eqnarray}
& & \myvec{u} = \left([T(s)]_1\right)_{i=1}^n - \myvec{t} \label{eq_uTt}\\
& & \myvec{v} = \left([s]_1\right)_{i=1}^n - \myvec{x} \label{eq_vsx}
\end{eqnarray}

Then from checks (1) and (2) in
$\glverify()$, one can infer the following, with
the prover's knowledge of
all randoms $r_u$, $r_u'$, $r_v$, $r_v'$:

\begin{small}
\begin{eqnarray}
\Z_u &=& \prod_{i=1}^n r^{2(i-1)} e(\myvec{u}_i, [1]_2)  ~e(r_u [1]_1, [1]_2) \label{zu} \\
\Commit_{\Z_u} &=& r_u [1] + r_u' [\nu]_1 \label{equ} \\
\Z_v &=& \prod_{i=1}^n r^{2(i-1)} e(\myvec{v}_i, \prf_i) 
	~e(r_v [1]_1, [1]_2) \label{rv} \\
\Commit_{\Z_v} &=& r_v [1] + r_v' [\nu]_1 \label{eqv}
\end{eqnarray}
\end{small}

They lead to:
\begin{small}
\begin{eqnarray}
\Z_u/\Z_v &=& \frac{
\prod_{i=1}^n r^{2(i-1)} e(\myvec{u}_i, [1]_2) }
{\prod_{i=1}^n r^{2(i-1)} e(\myvec{v}_i, \prf_i) }
~
\frac{e(r_u [1]_1, [1]_2) } {~e(r_v [1]_1, [1]_2) } \label{eq_dv1} \\
\Commit_{\Z_u} - \Commit_{\Z_v} & = &
	(r_u - r_v) [1] + (r_u'-r_v')[\nu]_1 \label{eq_dv2}
\end{eqnarray}
\end{small}

Then obvserve the $\prf_{\ttt{DL}}$'s check on
$\Z_u/\Z_v$ and $\Commit_{\Z_u} - \Commit_{\Z_v}$.
Note that the same $x$ is used in $\prf_{\ttt{DL}}$.
Note that $\Commit_{\Z_u}$ and $\Commit_{\Z_v}$ are
Pedersen commitment and they are computational binding.
The $\prf_{\ttt{DL}}$ establishes that
for an efficient prover there is one and only one
knowledge representation
of $\Commit_{\Z_u} - \Commit_{\Z_v} = x[1]_1 + y[\mu]_1$
where $x = r_u - r_v$ and $y = r_u' - r_v'$.
Since the $\prf_{\ttt{DL}}$ also proves that:
$\Z_u/\Z_v = e([1]_1, [1]_2)^{r_u-r_v}$, combining
with Equation\ \ref{eq_dv1},we have:
\begin{equation}
\prod_{i=1}^n r^{2(i-1)} e(\myvec{u}_i, [1]_2) 
= \prod_{i=1}^n r^{2(i-1)} e(\myvec{v}_i, \prf_i) 
\end{equation}

Given $r$ is the verifier challenge using
Fiat-Shamir transform, we immediately have:
\begin{equation}
\forall i\in [n]: 
	e(\myvec{u}_i, [1]_2) = e(\myvec{v}_i, \prf_i) \label{equv}
\end{equation}

Combinig Equations \ref{eq_uTt}, \ref{eq_vsx} and \ref{equv}, we then
have:
\begin{equation}
\forall \in in [n]:	
	e([T(s)]_1 - \myvec{t}_i, [1]_2) = e([s]_1 - [\myvec{x}_i]_1, 
		\prf_i) \label{eq_final}
\end{equation}

Given that each $\myvec{x}_i \in [n]$ (shown by the zk-CQ proof),
this completes the proof for knowledge soundness. For zero knowledge,
the proof is very similar to that of Section\ \ref{sec:eqcom}, i.e.,
all group elements in the proof are independent of each other (identical
to the distribution of real conversations between honest
prover and verifier). Then run simulators for each component
proof.
\end{proof}

It is possible to slightly improve the concrete preformance
of the protocol by replacing the first $\zktippprove$
by a $\zkmippprove$. We use the current implementation for its
convenience of supporting vector of $\G_2$ elements with
few changes in engineering efforts.
